\section{Bilan de la simulation}

\subsection{Points forts}

La simulation actuelle presente plusieurs qualites:

\begin{itemize}
  \item elle est structuree autour d'un paquet ROS 2 dedie;
  \item elle contient des mondes et scenarios reproductibles;
  \item elle expose les capteurs simules sur des topics ROS compatibles APEX;
  \item elle reutilise les memes noeuds de planification et suivi que la pile reelle;
  \item elle fournit une verite terrain et des outils de recording.
\end{itemize}

\subsection{Limites observees}

La complexite du launch \path{apex_sim.launch.py} est elevee. Ce fichier concentre beaucoup de modes, de parametres, de conditions et de chemins optionnels. Cette richesse est utile pour l'experimentation, mais elle augmente le cout de maintenance et la difficulte de lecture.

La coexistence de plusieurs simulations historiques peut aussi creer de la confusion: Gazebo Sim APEX, ancien \path{spawn_rc_car.launch.py}, simulation \path{voiture_system} avec Gazebo Classic, et ancien simulateur Python documente comme reference externe issue de \path{full_soft/Simulateur}.

\subsection{Interpretation}

D'apres l'analyse du depot, la simulation APEX est suffisamment riche pour servir de base principale aux essais. Elle doit cependant etre documentee et encadree pour eviter que les utilisateurs ne choisissent un ancien chemin par erreur.
